\chapter{Layout, pacotes e fontes}
\label{chap:layout-fontes}

\section{Estendendo LaTeX com pacotes}

Pacotes adicionam recursos ao documento. Eles são carregados no preâmbulo por
uma declaração própria. O pacote \arquivo{geometry}, por exemplo, configura
papel e margens:

\begin{lstlisting}[language=LaTeXBook]
\usepackage[a4paper,margin=2.5cm]{geometry}
\usepackage{graphicx} % inclusao de imagens
\usepackage{xcolor}  % definicao de cores
\usepackage{hyperref} % links no PDF
\end{lstlisting}

Carregue apenas o necessário e mantenha a configuração centralizada. Em um
livro com vários capítulos, um arquivo \arquivo{preamble.tex} evita que cada
capítulo redefina cores, margens e estilos.

\section{Por que LuaLaTeX?}

LuaLaTeX\index{LuaLaTeX} lê Unicode nativamente e, com \arquivo{fontspec}, usa fontes TrueType e
OpenType do sistema. Isso facilita criar uma identidade visual sem converter
arquivos de fonte para formatos antigos.

\begin{lstlisting}[language=LaTeXBook]
\usepackage{fontspec}
\setmainfont{TeX Gyre Pagella}
\setsansfont{TeX Gyre Heros}
\setmonofont{Latin Modern Mono}
\end{lstlisting}

A fonte principal atende ao texto corrido; a sem serifa é apropriada para
títulos e elementos gráficos; a monoespaçada atende código. Para usar uma fonte
instalada, substitua o nome, por exemplo:

\begin{lstlisting}[language=LaTeXBook]
\setmainfont{Libertinus Serif}
\setsansfont{Source Sans 3}
\setmonofont{JetBrains Mono}
\end{lstlisting}

Confira os nomes reconhecidos pelo sistema com \arquivo{fc-list} no Linux. Uma
mensagem ``font not found'' normalmente indica nome incorreto ou fonte ausente.

\section{Fontes armazenadas no projeto}

Guardar arquivos licenciados para redistribuição numa pasta \arquivo{fonts/}
torna o resultado reproduzível em outras máquinas:

\begin{lstlisting}[language=LaTeXBook]
\setmainfont[
  Path=fonts/,
  UprightFont=MinhaFonte-Regular.otf,
  BoldFont=MinhaFonte-Bold.otf,
  ItalicFont=MinhaFonte-Italic.otf
]{MinhaFonte}
\end{lstlisting}

Verifique sempre a licença. Algumas fontes permitem uso no PDF, mas não a
redistribuição do arquivo original.

\section{Cores e comandos próprios}

Nomes centralizados deixam alterações futuras mais seguras:

\begin{lstlisting}[language=LaTeXBook]
\definecolor{AzulLivro}{HTML}{183B56}
\newcommand{\termo}[1]{\textbf{\color{AzulLivro}#1}}
\end{lstlisting}

Depois, \verb|\termo{tipografia}| aplica sempre o mesmo tratamento. Defina
comandos pelo significado, não por uma aparência passageira. Um nome como
\comando{termo} é mais durável que um nome baseado em cor e tamanho.

\begin{pratica}
  Escolha uma fonte para texto, outra para títulos e uma monoespaçada. Altere
  uma cor e compile. Avalie legibilidade em parágrafos longos, itálico, negrito
  e código antes de confirmar a combinação.
\end{pratica}
